{
% \newcommand{\nVSLength}{0.170}
% \newcommand{\nVSEnd}{0.171}%slightly larger 0.001 than increment length to prevent breaking of chunks on PDF viewer.

\newlength{\nVSLength}	\setlength{\nVSLength}{0.56in}
\newlength{\nVSEnd}	\setlength{\nVSEnd}{0.001in+\nVSLength}
\newlength{\nRightSide}	\setlength{\nRightSide}{0in}
\newlength{\nLeftSide}	\setlength{\nLeftSide}{0in}

\psset{gradlines=100,
 fillstyle=gradient,
 gradangle=90,
 gradmidpoint=1.0,
 linewidth=0pt,
 linestyle=none,
 linearc=0pt}
{\Large Видимий спектр \hspace{0.07in}
	\psframebox[fillstyle=none,linestyle=none]{%
		\psset{xunit=2.8in}
		\multido{%
			\nCounter=0+1,
			\nStartColor=0.00+0.1,
			\nEndColor=0.1+0.1
			}{8}{%
			\setlength{\nRightSide}{\nVSEnd+\nVSLength*\nCounter}
			\setlength{\nLeftSide}{\nVSLength*\nCounter}
			\definecolor{StartColor}{hsb}{\nStartColor,1,1}
			\definecolor{EndColor}{hsb}{\nEndColor,1,1}
			\psframe[fillstyle=gradient,
				gradangle=90,
				gradbegin=StartColor,
				gradend=EndColor,
				gradmidpoint=1.0,
				linewidth=0pt,
				linestyle=none]
			(\nLeftSide,0in)(\nRightSide,0.15in)
		}
	}
}
\begin{itemize}

\item Діапазон EMR, видимий людині, називається ``світлом''. Видимий спектр також досить точно відповідає діапазону EMR, що проходить через нашу атмосферу від Сонця.

\item Інші істоти бачать різні діапазони видимого світла; наприклад, джмелі можуть бачити ультрафіолетове світло, а собаки мають іншу реакцію на кольори, ніж люди.

\item Небо блакитне, тому що наша атмосфера розсіює світло, і короткохвильовий блакитний колір розсіюється найбільше. Здається, що все небо освітлене блакитним світлом, але насправді це світло розсіяне від Сонця. Довші хвилі, такі як червоний та оранжевий, проходять прямо через атмосферу, що робить Сонце схожим на яскраву білу кулю, що містить усі кольори видимого спектра.

\item Цікаво, що видимий спектр охоплює приблизно одну октаву (подвоєння/половинення частоти).

\item Астрономи використовують фільтри для захоплення певних довжин хвиль та відсіювання небажаних. Основні астрономічні (візуальні) фільтри показані як  \psframebox[fillstyle=none,linestyle=none]{\astrofilter{0.07,0.03}{X}}

\end{itemize}
}
